% Generated by encode_cmp.py. Do not edit.
Each comparison table is three assembled strips. The columns are the ISA, the address, the encoding, and the instruction. RVV addresses are byte offsets in the acceptance image, and the hex is the instruction word. Intel bytes are x86-64 in memory order, with no spaces between bytes. They are assembled by clang with AVX2 and FMA, and without AVX-512, except \texttt{wide32} and \texttt{wide64}, whose Intel strip is EVEX because that register file is 512 bits. A splat from a general register is \texttt{movd} and a shuffle, or \texttt{vmovq} and \texttt{vpbroadcastq} from an \texttt{xmm}. ARM bytes are AArch64 NEON words, assembled by clang. SVE is not in these bytes. An indexed load on NEON is a scalar \texttt{ldr} and \texttt{str}. \texttt{psubusb} then \texttt{pcmpeqb} is the unsigned compare. \texttt{pcmpgtb} is signed and is not used. \texttt{pmullw} alone is the low half of a product. The \texttt{vwmul16} Intel strip is \texttt{pmullw}, \texttt{pmulhw}, unpack, and \texttt{packssdw}. The ARM strip is \texttt{smull} and \texttt{smull2}, then \texttt{sqxtn} and \texttt{sqxtn2}. \texttt{vfmadd231ps} and AArch64 \texttt{fmla} are fused. A separate multiply and add is not. The GEMM strip is one rank-1 update: a vector load, a scalar broadcast, and a fused multiply-add into each accumulator. \texttt{vbroadcastss}, \texttt{vmovddup}, and \texttt{ld1r} are that broadcast. \texttt{wide32} and \texttt{wide64} are that update on a 512-bit register, four columns, stride \texttt{vl}. A 24 by 8 tile is still wider than those four accumulators. The ARM column of that strip stays NEON. A label shares its row with the following instruction. On \texttt{asum} and \texttt{iamax} the Intel mask in \texttt{xmm7} is the absolute-value mask, the same role as the byte mask in \texttt{mask8}. The index itself is the RVV \texttt{vfirst.m}. \texttt{quad}, \texttt{add128}, \texttt{addd16}, \texttt{zext64}, and \texttt{sh64} are RV128I integer strips. Their first column is labeled RV128. \texttt{lq}, \texttt{sq}, \texttt{ldu}, and \texttt{addid} use the decoder placement in Chapter~\ref{ch:rv128}. \texttt{addi} is the ordinary XLEN add. \texttt{slli} by 64 uses the seven-bit immediate. The Intel and ARM columns of those strips keep 64-bit general registers and write the high half explicitly, except \texttt{quad}, which moves 16 bytes through an \texttt{xmm} and through NEON. SVE and AVX-512 are not in those bytes. \texttt{sad16} is a halfword absolute difference and a horizontal sum. \texttt{pabsw} is the SSSE3 absolute value. \texttt{psadbw} sums bytes into words and is a different instruction. \texttt{sw8} is \texttt{paddusb} and \texttt{pmaxub}. \texttt{wmul32}, \texttt{bfly32}, and \texttt{cmul32} are \texttt{vmulps}, \texttt{vaddps}, and \texttt{vsubps}. Those three are not fused.
